\chapter{Requirements Specification \& Documentation (SRS)}
\section{Functional Requirements}
\subsection{Conventions and System Boundary}
Table~\ref{tab:functional} restates each requirement in plain English as the required behavior plus an acceptance interpretation. Workbook identifiers are retained for traceability; the original wording, complexity, sources, dependencies, and statuses are in \texttt{data/dataset.json}. Proposed specification corrections are identified separately and do not modify the workbook.

External actors provide plate numbers, sensor readings, staff commands, or fire-alarm signals. The simulation manages slots, tickets, and transactions; opening a barrier changes a software state variable. Successful business workflows do not establish hardware accuracy or resistance to real-world fraud.
\clearpage
\par\begingroup\footnotesize\setlength{\tabcolsep}{4pt}\renewcommand{\arraystretch}{1.1}
\begin{longtable}{>{\raggedright\arraybackslash}p{\dimexpr 0.22000\linewidth-2\tabcolsep\relax}>{\raggedright\arraybackslash}p{\dimexpr 0.32000\linewidth-2\tabcolsep\relax}>{\raggedright\arraybackslash}p{\dimexpr 0.46000\linewidth-2\tabcolsep\relax}}
\caption{Functional specifications and verification criteria}\label{tab:functional}\\
\toprule
\textbf{ID / name} & \textbf{Required behavior} & \textbf{Acceptance / verification interpretation} \\ \midrule
\endfirsthead
\multicolumn{3}{l}{\small\itshape Functional specifications and verification criteria (continued)}\\
\toprule
\textbf{ID / name} & \textbf{Required behavior} & \textbf{Acceptance / verification interpretation} \\ \midrule
\endhead
\midrule
\multicolumn{3}{r}{\scriptsize Continued on the next page}\\
\endfoot
\bottomrule\endlastfoot
REQ-01 Vehicle entry & Issue a ticket with ID, entry time, and a matching slot when capacity exists. & Create the ticket; mark the slot Occupied; reject a plate already in the lot. \\
REQ-02 Availability display & Show free-slot counts by floor and zone. & Displayed counts equal Free leaves in each branch. \\
REQ-03 Exit and fee calculation & Calculate fees using started hours and the vehicle-type rate. & Outside peak hours, without surcharge or exemption, fee equals rounded-up hours times the rate. \\
REQ-04 Payment & Confirm payment before opening the exit barrier in normal mode. & Failed payment retains the vehicle and ticket; allow another method. Evacuation takes priority. \\
REQ-05 Reservation & Hold a matching slot until 15 minutes after expected arrival. & Do not assign it to another vehicle; release a no-show only after the hold expires. \\
REQ-06 Locate a vehicle & Look up a parked vehicle by plate. & Return floor, zone, slot, and distance; report not found without an active ticket. \\
REQ-07 Lost ticket & Look up the plate, verify ownership, charge the fee, and record the incident. & No acceptance criteria in the workbook; proposed criteria and fee discrepancies are discussed in Section 3.1.3. \\
REQ-08 Zone lock/unlock & Allow administrators to lock or unlock a floor or zone. & Locked slots are not assigned to new vehicles; existing vehicles remain and exit through the appropriate workflow. \\
REQ-09 Access control & Restrict administrative functions to authenticated roles. & Attendants cannot lock zones; only administrators change configuration. \\
REQ-10 Transaction audit log & Record entry, exit, and payment transactions. & Events carry timestamps and can be retrieved for a reporting period. \\
REQ-11 Plate recognition & Accept recognition confidence of at least 80\%; otherwise require manual input. & Accept 80\%; 79\% or camera failure requires manual confirmation. \\
REQ-12 Nearest-slot allocation & Select the nearest Free slot of the correct type. & Exclude locked/faulty slots; respect reservations and VIP eligibility; break distance ties by tree order. \\
REQ-13 Sensor reconciliation & Alert when sensor readings differ from ticket data. & Report unticketed or missing vehicles; a silent sensor makes its slot Blocked. \\
REQ-14 Overstay alert & Alert when parking exceeds 24 hours. & A sensor scan detects and records a stay longer than 24 hours. \\
REQ-15 Emergency evacuation & Open every barrier, stop new tickets, and waive exit fees. & Triggered by fire alarm or administrator; cancel reservations; only an administrator may end evacuation. \\
REQ-16 Peak-hour pricing & Apply a 1.5 multiplier to hours starting within 07-09 or 17-19. & Evaluate each started hour in UTC+7; the upper boundary is excluded. \\
REQ-17 Occupancy surcharge & Apply 25\% when same-type occupancy exceeds 80\% at issuance. & Store the multiplier on the ticket; code checks after allocation and includes Reserved slots. \\
REQ-18 Staff authentication & Authenticate with a password and lock after three consecutive failures. & Valid login returns a role; the third failure locks the account in the current Accounts object. \\
REQ-19 Cancel reservation & Cancel a reservation and release its slot immediately. & The slot becomes available or Blocked if locked; distinguish expired reservations. \\
REQ-20 Lot/rate configuration & Allow administrators to change rates and add valid zones. & New rates affect only new tickets; reject nonpositive rates and duplicate zone names. \\
REQ-21 Revenue report & Aggregate revenue and audit records over a period. & Zone totals reconcile with floor and overall totals; access is restricted to manager/admin. \\
REQ-22 Fee override & Record a reasoned request for manager/admin approval. & Amount is nonnegative; attendants cannot approve; an approved amount replaces the exit total. \\
REQ-23 Monthly/VIP passes & Manage plate-based passes and waive parking for valid passes. & Expired passes at entry use normal pricing; VIP slots require VIP eligibility; energy remains chargeable. \\
REQ-24 EV slots & Allocate dedicated slots to electric vehicles. & EVs do not receive car slots and vice versa; reject EV entry when EV capacity is exhausted. \\
REQ-25 Charging fee & Add total kWh times 3,500 VND to the fee. & Accept positive kWh only for EV tickets; two normal hours and 10 kWh total 59,000 VND. \\
REQ-26 Exit-plate verification & Compare the exit-camera plate with the ticket plate. & Mismatch or unreadable plates block exit and are logged; authorized confirmation may release the vehicle. \\
REQ-27 Grace period & Waive parking fees for stays of at most 15 minutes. & 15 minutes costs zero; 16 minutes costs one started hour; do not deduct the grace period from longer stays. \\
REQ-28 Daily fee cap & Cap each 24-hour block at ten hourly rates before surcharge. & 24 normal hours cost 100,000 VND for a car; 25 normal hours cost 110,000 VND. \\
REQ-29 Offline operation & Continue local entry/exit when disconnected. & Simulated cash/card and evacuation work; refuse online payment, reservations, and pass changes. \\
REQ-30 Recovery synchronization & Mark offline records and process them on reconnection. & Clear offline marks and log the count; the prototype does not transmit records to a real server. \\
REQ-31 Barrier interlock & Refuse to close a barrier while a vehicle is detected or the presence sensor fails. & A close command with a vehicle present, a sensor fault, or power loss leaves the barrier open and is logged; a clear way allows closing. \\
\end{longtable}\endgroup


\subsection{Pricing Rules and Order of Application}
Sample rates per started hour are 10,000 VND for cars, 3,000 VND for motorcycles, and 12,000 VND for EVs. Charging costs 3,500 VND/kWh. Time is evaluated in UTC+7. Each started hour is anchored to the entry time rather than aligned to clock-hour boundaries.

Let $\Delta t$ denote the stay in hours, $r$ the rate recorded on the ticket, $n=\lceil\Delta t\rceil$, and $m_i$ a multiplier of 1.5 when the start of hour $i$ falls within [07:00,09:00) or [17:00,19:00), and 1 otherwise. For a stay longer than 15 minutes without a monthly/VIP pass, the simulated parking fee is:
\begin{equation}
F_{\mathrm{parking}}=\left\lfloor s\sum_{d=0}^{\lceil n/24\rceil-1}\min\left(10r,\sum_{i=24d}^{\min(24d+23,n-1)} r m_i\right)\right\rfloor .
\end{equation}
Here $s=1.25$ when occupied or reserved slots of the vehicle type exceed 80\% at ticket issuance; otherwise $s=1$. The current code calculates occupancy after marking the arriving vehicle's slot Occupied. The $10r$ cap is applied before the surcharge, in 24-hour blocks measured from entry. Consequently, the cap after surcharge can reach $12.5r$. This ordering requires explicit agreement when approving REQ-17 and REQ-28.

The parking component is zero for stays of at most 15 minutes or for a valid monthly/VIP pass recorded at entry. Energy charges still apply. The normal total is:
\begin{equation}
F_{\mathrm{total}}=F_{\mathrm{parking}}+\operatorname{round}(3500E),
\end{equation}
where $E$ is the total energy delivered in kWh. An approved override replaces the entire total. During evacuation, exit costs zero and bypasses payment checks. The grace period is not deducted from a stay once the stay exceeds 15 minutes.

For example, a car entering at 16:20 and leaving at 19:05 has three started hours at 16:20, 17:20, and 18:20. They cost 10,000, 15,000, and 15,000 VND, giving 40,000 VND without an occupancy surcharge. An EV parked for two off-peak hours and charged with 10 kWh costs 59,000 VND. These examples can be checked against \texttt{live.py} and TC-58.

\subsection{Specification Issues to Resolve}
REQ-07 lacks acceptance criteria and combines lookup, charging, logging, notification, and receipt printing. It should be split into child requirements. Minimum acceptance conditions should cover finding the ticket by plate, refusing exit without ownership evidence, recording a timestamped incident, and collecting the approved fee. The code currently charges $24r+50{,}000$ VND for a lost ticket, while REQ-28 defines a daily cap of $10r$ before surcharge. The business owner must clarify the phrase ``maximum daily fee'' in REQ-07 before acceptance.

The tool flags REQ-10 as potentially non-atomic using a string-based heuristic. Logging could be normalized into three event-specific requirements beneath one parent. In addition, exit can skip plate checking when \texttt{seen} is \texttt{None}. Missing input must be distinguished from a camera explicitly returning an unreadable plate. Tests of the mismatch branch do not establish that every call requires plate verification.

\section{Non-functional Requirements}
\par\needspace{11\baselineskip}\begingroup\footnotesize\setlength{\tabcolsep}{4pt}
\begin{longtable}{>{\raggedright\arraybackslash}p{\dimexpr 0.23000\linewidth-2\tabcolsep\relax}>{\raggedright\arraybackslash}p{\dimexpr 0.52000\linewidth-2\tabcolsep\relax}>{\raggedright\arraybackslash}p{\dimexpr 0.25000\linewidth-2\tabcolsep\relax}}
\caption{Non-functional requirements}\label{tab:nfr-spec}\\
\toprule
\textbf{ID / category} & \textbf{Specification and target} & \textbf{Priority / verification} \\ \midrule
\endfirsthead
\multicolumn{3}{l}{\small\itshape Non-functional requirements (continued)}\\
\toprule
\textbf{ID / category} & \textbf{Specification and target} & \textbf{Priority / verification} \\ \midrule
\endhead
\midrule
\multicolumn{3}{r}{\scriptsize Continued on the next page}\\
\endfoot
\bottomrule\endlastfoot
NFR-01 / Performance & Slot-assignment p95 at most one second. & Must / Test \\
NFR-02 / Availability & Availability during operating hours at least 99.5\%. & Must / Analysis \\
NFR-03 / Usability & New attendants handle a lost ticket within two minutes after 20 minutes of training. & Should / Demonstration \\
NFR-04 / Security & Encrypt 100\% of PII fields at rest and in transit; target AES-256 / TLS 1.2+. & Must / Inspection \\
NFR-05 / Scalability & Support growth in floors and slots; metric and target are missing. & Could / Analysis \\
NFR-06 / Safety & All barriers open within five seconds of an evacuation trigger. & Must / Test \\
\end{longtable}\endgroup

These targets describe required system properties. NFR-04 encryption is not implemented in the prototype: demonstration passwords are plaintext in memory, workbook and JSON files are unencrypted, and no actual TLS channel exists. NFR-05 lacks both a metric and a threshold, preventing acceptance. Chapter 6 separates measurements, simulation analysis, and unverified targets.

\section{Chapter Conclusion}
The SRS covers the main business capabilities and gives verification conditions for most requirements. Pricing, evacuation priority, and data-state rules must be agreed before real deployment. The REQ-07 and NFR-05 defects remain visible so that baseline quality is reported transparently.

Three decisions need attention. The fee rules apply the daily cap before the occupancy surcharge, so the effective cap can reach $12.5r$ rather than $10r$; the business owner must approve this ordering. The lost-ticket rule ($24r+50{,}000$ VND) must be reconciled with the daily cap in REQ-28 and with REQ-07. Finally, NFR-04 encryption is specified but not implemented in the prototype, and NFR-05 has neither metric nor threshold. These items return in Chapters~6 and~7 as verification limits and as quality findings.
